\begin{textsample}{POS Dim 3 – persona\_gemini – Score 17.00 – t551\_gemini.txt}  \label{ex:f3_pos_049}
The \textbf{car} \textbf{industry} is pursuing a dangerous \textbf{trend} of producing increasingly heavier vehicles.
This push towards larger \textbf{cars}, particularly SUVs, \textbf{means} they consume more \textbf{energy} and, as a result, emit higher \textbf{levels} of CO2.
This \textbf{problem} is not solved by electrification, as heavy electric \textbf{models} require larger, resource-intensive batteries, continuing the \textbf{cycle} of high \textbf{energy} \textbf{consumption}.
A clear example of this \textbf{trend} is the Volkswagen Golf, which has gained significant weight over the decades, resulting in a 57 % \textbf{increase} in its \textbf{emissions}.
The \textbf{company} 's heavier Tiguan \textbf{model} further illustrates this issue.
Beyond the environmental impact, these oversized vehicles create significant safety hazards.
SUVs pose a greater risk in accidents, especially to pedestrians and other road users.
Furthermore, their \textbf{drivers} are noted to be more prone to engaging in risky behaviour.
Faced with this reality, the \textbf{car} \textbf{industry} continues to prioritise profits over its responsibilities under the Paris Agreement.
It invests heavily in SUVs, promotes them with extensive advertising, and has seen \textbf{sales} boom as a result.
A sustainable \textbf{transport} future requires a different path.
We \textbf{need} a transition to smaller, lighter electric \textbf{cars}, an \textbf{increase} in shared ownership \textbf{schemes}, and a major expansion of \textbf{public} \textbf{transport}.
The ultimate goal should be fewer vehicles on our roads.
It is this vision that has prompted Greenpeace to take direct action, protesting against the import of oversized SUVs.

% matched lemmas: car, company, consumption, cycle, driver, emission, energy, increase, industry, level, mean, model, need, problem, public, sale, scheme, transport, trend
\end{textsample}
